\subsection{相机参数}
本小节将会结合相机的结构去介绍相机的各种参数。
\subsubsection{感光元件}
相机目前的主流的感光元件有CCD（电荷耦合）和CMOS（互补金属氧化物半导体）两种。

\textbf{CCD}

Charge Coupled Device，它使用一种高感光度的半导体材料制成，由许多感光单位组成，通常以百万像素为单位。当CCD表面受到光线照射时，每个感光单位会将电荷反映在组件上，即把光线转变成电荷，电荷在像素之间依次传递，最终集中到串联寄存器，转换为电压信号，然后放大，经模数转换后形成数字信号。而后数字信号，经过压缩后保存在相机内部的闪速存储器或内置硬盘卡中。有能力生产CCD 的公司分别为：索尼、飞利浦、柯达、松下、富士、夏普，大半是日本厂商。

\textbf{CMOS}

Complementary Metal-Oxide Semiconductor，和CCD一样同为在数码相机中可记录光线变化的半导体。CMOS的制造技术和一般计算机芯片没什么差别，主要是利用硅和锗这两种元素所做成的半导体，使其在CMOS上共存着带N（带–电） 和 P（带+电）级的半导体，这两个互补效应所产生的电流即可被处理芯片纪录和解读成影像。
每个像元都有独立的节点，转换为电压信号，然后放大，经模数转换后形成数字信号。

\textbf{CCD与CMOS的区别}

1.成像过程

CCD仅有一个(或少数几个)输出节点统一读出，其信号输出的一致性非常好。

CMOS中每个像素都有各自的信号放大器，各自进行电荷-电压转换，信号输出的一致性较差。

CCD为了读出整幅图像信号，要求输出放大器的信号带宽较宽。

CMOS中，每个像元中的放大器的带宽要求较低，大大降低了芯片的功耗，但数以百万的放大器的不一致性却带来了更高的固定噪声。

2.集成性

CCD中的电路和器件是集成在半导体单晶材料上，工艺较复杂。其仅能输出模拟电信号，需要后续的地址译码器、模数转换器、图像信号处理器处理，并且还需要提供三组不同电压的电源和同步时钟控制电路，集成度非常低。

CMOS能将图像信号放大器、信号读取电路、A/D转换电路、图像信号处理器及控制器等集成到一块芯片上，集成度高，芯片级相机的概念也由此产生。

3.速度

CCD采用光敏元输出，只能按规定的程序输出，速度较慢。

CMOS有多个电荷-电压转换器和行列开关控制，读出速度较快。此外，CMOS的地址选通开关可随机采样，实现子窗口输出，在仅输出子窗口图像时可获得更高的速度。

4.噪声

CCD采用PN结或二氧化硅隔离层隔离噪声，成像质量相对CMOS有一定的优势。

CMOS的集成度较高，各元件、电路之间距离很近，干扰较为严重，噪声对图片质量影响很大。

5.功耗

CCD需要3路电源来满足特殊时钟的需求，功耗较大。

CMOS只需一个电源供电，功耗仅为CCD的1/10。

\textbf{总结}

目前，市场销售的数码摄像头中以CMOS感光器件的为主。在采用CMOS为感光元器件的产品中，通过采用影像光源自动增益补强技术，自动亮度、白平衡控制技术，色饱和度、对比度、边缘增强以及伽马矫正等先进的影像控制技术，完全可以达到与CCD摄像头相媲美的效果。

受市场情况及市场发展等情况的限制，摄像头采用CCD图像传感器的厂商为数不多，主要原因是采用CCD图像传感器成本高的影响，同时由于CCD耗电量大，工艺复杂，像素提高难度加大。

\subsubsection{相机感光参数}
调整相机对环境的感光程度，主要有以下三个途径，分别是光圈、快门和ISO。

\textbf{光圈}

通过调整镜头光圈的大小，调整进光量，光圈大小有两种写法：大F写法（F2.8）、小f写法（f/2.8），两种写法完全等价，例如F4指的是光圈大小占镜头直径的4分之1。
F后面的数字越大，光圈越小，反之则反。

1.挡位设计

光圈的档位设计是相邻的两档的数值相差1.4倍（2的平方根1.414的近似值）相邻的两档之间。
透光孔直径相差根号2倍，透光孔的面积相差一倍， 底片上形成的影像的亮度相差一倍，维持相同曝光量所需要的时间相差一倍。

所有的光圈挡位有：

f/1.0，f/1.4，f/2.0，f/2.8，f/4.0，f/5.6，f/8.0，f/11，f/16，f/22，f/32，f/44，f/64

2.光圈分类

固定光圈：最简单的相机只有一个园孔的固定光圈

沃特侯瑟光圈：最初的可变光圈只是一系列大小不同的园孔排列在一个有中心轴的圆盘的周围；转动园盘可将适当大小的圆孔移到光轴上,达到控制孔径的效果。十九世记中叶约翰·沃特侯瑟发明这种光圈。

猫眼式光圈：由一片中心有椭圆形或菱形孔的金属薄片平分为二组成，将两片有半椭圆形或半菱形孔的金属薄片对排,相对移动便可形成猫眼式光圈。猫眼式光圈多用于简单照相机。

"虹膜"式光圈：是由多个相互重叠的弧形薄金属叶片组成的，叶片的离合能够改变中心圆形孔径的大小。有些照相机可以借助转动镜头筒上的圆环改变光圈孔径的大小，而有些照相机则是利用微处理器芯片控制微电机自动地改变光圈的孔径。弧形薄金属叶片可多达十八片。弧形薄金属叶片越多，孔径越近圆形。通过电子计算机设计薄金属叶片的形状，可以只用7片薄金属叶，得到近圆形孔径。

瞬时光圈：单反照相机的光圈是瞬时光圈，只在快门开启的瞬间，光圈缩小到预定大小。平时光圈在最大位置。

兼快门光圈：有的简便照相机的光圈兼有快门的功能，这类兼快门光圈大多是双叶片的猫眼式光圈，与单纯猫眼式光圈不同的是：于兼快门光圈平时是完全关闭的：在按下快门的瞬间，双叶片光圈开启到预定的孔径后，保持这孔径到一段预定快门开启时间之后，立刻闭合：如此一来，光圈便又兼快门的功能。

偷懒不放图了先（逃

\textbf{快门}

快门是遮光板，只有当快门打开的时候，底片才能接收到光照。所以快门可以控制进光时间。
快门越块：成像时间越短，被拍摄物品移动的距离越小，图像越清晰。常用于静态拍摄。
反之则反

但是随着相机技术进步，现在相机大部分使用电子快门，常见的电子快门主要有两种，卷帘快门和全局快门

1.卷帘快门

卷帘快门通过逐行扫描的方式从上到下捕获图像，而不是一次性捕获所有像素。这意味着每一行像素的曝光存在微小的时间差。当相机或物体高速移动时，这种逐行扫描可能导致图像失真，形成“果冻效应”
但是其有着成本低，分辨率高的优点，使其在一些需要高分辨率但对动态性能要求不高的应用中非常适用，如静态图像捕捉或家用监控。

卷帘快门导致的果冻效应，其显著程度主要取决于以下几个因素：

*物体的移动速度：移动速度越快，失真越明显。

*物体的大小：大尺寸物体在画面中的扭曲更加明显。

*分辨率：更高的分辨率在相同帧率下行时间更短，相对可以减少失真。

2.全局快门

全局快门（Global Shutter）是一种图像传感器技术，其核心在于能够同时捕获画面中所有像素的信息。与逐行扫描的滚动快门不同，全局快门的所有像素在同一时间曝光，然后一起读取数据，这种同步曝光方式使其在高速运动场景中能够有效消除图像失真现象，如“果冻效应”（rolling shutter effect）。可以将其比作拍照时快门在一瞬间“拍下”整个画面，而非逐行扫描捕捉，确保每一部分都与真实场景一致。

总结

静态场景可以使用卷帘快门相机，相同价格下，可以获得更高分辨率的相机，动态场景下，则需要根据情况权衡，优先全局快门相机，其次高帧率，高分辨率的
卷帘快门，在要求不高的程度上，也可以，请根据成本酌情考虑。在笔者眼里，全局快门的缺点就是贵（哭

\textbf{ISO}

感光度，又称为ISO值，是衡量底片对于光的灵敏程度。
对光越灵敏，越适合捕捉暗光线：亮处用小ISO即可，暗处用大ISO。对光越灵敏，意味着所需的进光时间越短：
意味着所需快门的时间越短。所以ISO值上限决定了快门的上限，只有ISO足够大，快门才可以足够快。
意味着曝光量越少。所以ISO越大，图像越模糊，使用较高敏感度通常会导致影像质量降低（较粗的底片颗粒或是较高的影像噪声）。
ISO每差一倍，曝光量差一档，也就是说相差一档光圈或者一档快门。

提高ISO会降低图像质量，所以优先采取其他方法提高亮度。

\subsubsection{图片噪点}
噪点在图像上常表现为一引起较强视觉效果的孤立像素点或像素块。比如上午所提的提高ISO或者相机长时间工作过热等，都可能会导致出现噪点。

图像常见噪声基本上有以下四种，高斯噪声，泊松噪声，乘性噪声，椒盐噪声。

\textbf{高斯噪声}

高斯噪声是指它的概率密度函数服从高斯分布（即正态分布）的一类噪声。如果一个噪声，它的幅度分布服从高斯分布，而它的功率谱密度又是均匀分布的，则称它为高斯白噪声。高斯白噪声的二阶矩不相关，一阶矩为常数，是指先后信号在时间上的相关性。

产生原因：

1）图像传感器在拍摄时市场不够明亮、亮度不够均匀。

2）电路各元器件自身噪声和相互影响。

3）图像传感器长期工作，温度过高。

\textbf{泊松噪声}

泊松噪声，就是符合泊松分布的噪声模型，泊松分布适合于描述单位时间内随机事件发生的次数的概率分布。如某一服务设施在一定时间内受到的服务请求的次数，电话交换机接到呼叫的次数、汽车站台的候客人数、机器出现的故障数、自然灾害发生的次数、DNA序列的变异数、放射性原子核的衰变数等。

\textbf{乘性噪声}

乘性噪声一般由信道不理想引起，它们与信号的关系是正相关，信号在它在，信号不在他也就不在。

\textbf{椒盐噪声}

椒盐噪声，椒盐噪声又称脉冲噪声，它随机改变一些像素值，是由图像传感器，传输信道，解码处理等产生的黑白相间的亮暗点噪声。


进度到这里，我是摄影高手（￣︶￣）↗　



接着由iso引出噪点介绍高斯/泊松/椒盐噪声模型

相机伽马值矫正

几种图片储存格式 贝尔曼/霍夫曼/YUV/RGB/HSV/HSL/CMYK/CMY/YCbCr/XYZ/LAB/LCH/LUV/HCL

景深是什么

相机对焦到无限远的原理

相机对于焦段的选择

相机镜头的选择

镜头参数

光圈大小之类的


\subsection{相机畸变模型}
\subsubsection{针孔相机}
\subsubsection{鱼眼相机}
\subsubsection{广角相机}

相机标定介绍
几种常用的标定板/棋盘,圆点,aruco,charuco

ros包标定 opencv标定 手写标定
每种给出文档，手写标定研究一下，写出来给

单目标定，双目标定，动态标定，在线标定，自标定（这两个有区别），无标定板标定等多种类型的标定

\subsubsection{相机选购}
几种常见的相机像素单位，商家说高清，不一定真高清

相机帧率不一定是当前的分辨率，可能指的是120hz的时候480p，1080p的时候其实只有30hz

分清黑白，分辨率等参数

定焦，自动调焦AF/FF
\subsection{相机安装}

efcedsf
edfewds
ewfews

ewfewrdfewrf

\subsubsection{786}

\textbf{加粗}

\begin{tcode}
#include <iostream>
using namespace std;

int main() {
    cout << "Hello, World!" << endl;
    return 0;
}
\end{tcode}

\begin{center}

        \begin{tabular}{ccccc}
            \hline
            捕获方式 & 拷贝方式 & 是否可以修改外部值 & 语法 & 全捕获\\
            \hline
            值捕获 & 深拷贝 & 不可以 & [a,b] & [=]\\
            引用捕获 & 浅拷贝 & mutable，可以修改 & [\&a,\&b] & [\&]\\
            \hline
            
        \end{tabular}
    \end{center}

    \begin{enumerate}
        \item 1312
        \item 13131
        \item 21423423
        \item 2142134
    \end{enumerate}


    $int$

    $\binom{1}{2} $

    \# \_